\section{Proof of the exact tower law}\label{app:tower}
\subsection{Logarithmic rank and a quantitative transfer inequality}

For any probability density $p$ on $(0,\infty)$, set
\[
 g_p(y)=e^{y/2}\sqrt{p(e^y)},\qquad y\in\mathbb R,
 \qquad (T_hg)(y)=g(y-h).
\]
Then $\norm{g_p}_2=1$ and, for $h=\log d$,
\begin{equation}\label{eq:exact-rank-translation}
 \mathcal E_d(p):=
 \int_0^\infty
 \left(\sqrt{p(x)}-d^{-1/2}\sqrt{p(x/d)}\right)^2dx
 =\norm{g_p-T_hg_p}_2^2.
\end{equation}
When $p(x)=p_j$ on $(j-1,j]$ and $d$ is an integer, the integral is
exactly the discrete Schmidt-coefficient error.  In particular, the
passage to a continuous rank variable introduces no approximation.

\begin{lemma}[Hellinger transfer]\label{lem:hellinger-transfer}
For two probability densities $p,q$ on $(0,\infty)$ and any $d>0$,
\begin{equation}\label{eq:hellinger-transfer}
 \left|\sqrt{\mathcal E_d(p)}-
             \sqrt{\mathcal E_d(q)}\right|
 \leq 2\norm{\sqrt p-\sqrt q}_2.
\end{equation}
Consequently, if $b_n\to\infty$,
$\mathcal E_d(q_n)=a/b_n+O(b_n^{-3/2})$, and
$\norm{\sqrt{p_n}-\sqrt{q_n}}_2=O(b_n^{-1})$, then
$\mathcal E_d(p_n)=a/b_n+O(b_n^{-3/2})$.
If the latter distance is only $o(b_n^{-1/2})$, the conclusion is
$\mathcal E_d(p_n)=a/b_n+o(b_n^{-1})$.
\end{lemma}

\begin{proof}
The change of variables $x=e^y$ is an isometry on square roots of
densities.  The operator $I-T_h$ has norm at most two.  Apply the reverse
triangle inequality to $(I-T_h)g_p$ and $(I-T_h)g_q$ and then square.
\end{proof}

\subsection{The smooth Calabrese--Lefevre distribution}

For $b>0$, define $q_b$ on $(0,\infty)$ by
\begin{equation}\label{eq:CL-density-exact}
 q_b(x)=
 \begin{cases}
 e^{-b},&0<x\leq1,\\
 \displaystyle e^{-b-z(x)^2/(4b)},&x>1,
 \end{cases}
 \qquad I_0(z(x))=x,\quad z(x)>0.
\end{equation}
The constant piece $(0,1]$ represents the largest eigenvalue.  It is
essential for exact normalization; it has exponentially small weight
in the estimates below.

\begin{proposition}[Quantitative smooth law]\label{prop:CL-sharp-law}
The function $q_b$ is a probability density.  For every fixed $d>0$,
as $b\to\infty$,
\begin{equation}\label{eq:CL-law-rigorous}
 -\int_0^\infty q_b(x)\log q_b(x)\,dx=2b,
 \qquad
 \mathcal E_d(q_b)=\frac{(\log d)^2}{8b}
                   +O_d(b^{-3/2}).
\end{equation}
The logarithmic-rank square root $g_b=g_{q_b}$ obeys
\begin{equation}\label{eq:CL-derivative-bounds}
 \norm{g_b'}_2^2=\frac1{8b}+O(b^{-3/2}),
 \qquad
 \norm{(g_b'')_{\rm ac}}_2=O(b^{-1}).
\end{equation}
The only singular part of $g_b''$ is the atom
$-e^{-b/2}(2b)^{-1}\delta_0$.
\end{proposition}

\begin{proof}
Write $t=z^2/(4b)$.  The measure induced by $q_b(x)\,dx$ on $t\geq0$
has an atom $e^{-b}$ at zero and continuous part
\[
 e^{-b-t}\,dI_0(2\sqrt{bt}).
\]
Using $I_0(2\sqrt{bt})=\sum_{k\geq0}(bt)^k/(k!)^2$, Tonelli's theorem
gives, for $s>-1$,
\begin{equation}\label{eq:CL-compound-Poisson}
 \mathbb E e^{-sT}
   =e^{-b}\left(1+\int_0^\infty
       e^{-(1+s)t}\,dI_0(2\sqrt{bt})\right)
   =\exp\!\left(\frac{b}{1+s}-b\right).
\end{equation}
At $s=0$ this proves normalization.  Equivalently, $T$ is the sum of
a Poisson($b$) number of independent exponential random variables of
mean one.  Its cumulants are $\kappa_k(T)=b\,k!$, so
\begin{equation}\label{eq:CL-moments}
 \mathbb ET=b,\qquad
 \mathbb E(T-b)^2=2b,\qquad
 \mathbb E(T-b)^4=12b^2+24b.
\end{equation}
Since $-\log q_b=b+T$, its entropy is exactly $2b$.

For $y>0$, put $z=z(e^y)$ and
\[
 R(z)=I_1(z)/I_0(z),\qquad B(z)=z/R(z),\qquad
 a_b(y)=\frac{d}{dy}\log\bigl(e^yq_b(e^y)\bigr)
        =1-\frac{B(z)}{2b}.
\]
For $y<0$ the score $a_b$ is one.  The following elementary bounds
are used:
\begin{equation}\label{eq:bessel-ratio-bounds}
 z\leq B(z)\leq z+2\quad(z>0),\qquad
 \sup_{z>0}\left|\frac{B'(z)}{R(z)}\right|<\infty.
\end{equation}
For completeness, $0<R<1$ follows from the integral representation
of $I_0$ and $I_1$.  The Riccati equation
$R'=1-R^2-R/z$, together with the expansions at zero, shows that
$R(z)>z/(z+2)$: at any prospective crossing of that lower comparison
function, the derivative of their difference is
$3z/(z+2)^2>0$.  This proves the first pair of bounds.
At zero, $R=z/2+O(z^3)$ and $B=2+z^2/4+O(z^4)$, whence
$B'/R\to1$.  At infinity the fixed-order Bessel expansions give
$R=1-1/(2z)+O(z^{-2})$ and $B'=1+O(z^{-2})$.
Thus $B'/R\to1$ there as well, and continuity on compact subintervals
proves the last bound.  The recurrence and differentiated expansions
used here are also recorded in NIST DLMF, Sections 10.29 and 10.40.

Under the probability measure in \eqref{eq:CL-compound-Poisson}, let
$Z=2\sqrt{bT}$ and $U=T-b$.  The exact identity
\[
 Z-2b=U-\frac{U^2}{(\sqrt T+\sqrt b)^2}
\]
and \eqref{eq:CL-moments} imply
\begin{equation}\label{eq:CL-radial-moments}
 \norm{Z-2b-U}_{L^2}=O(1),\qquad
 \mathbb E(Z-2b)^2=2b+O(\sqrt b),\qquad
 \mathbb E|Z-2b|^4=O(b^2)
 \quad(b\geq1).
\end{equation}
For the last estimate use $|Z-2b|\leq2|U|$.
Extend $B$ continuously by $B(0)=2$.  Except at the atom $T=0$,
the score is $1-B(Z)/(2b)$.  At that atom its actual value is one;
substituting $1-B(0)/(2b)$ changes its second and fourth moments by
exponentially small quantities.  Therefore
\eqref{eq:bessel-ratio-bounds} and \eqref{eq:CL-radial-moments} give
\begin{equation}\label{eq:CL-score-estimates}
 \int a_b^2g_b^2\,dy=\frac1{2b}+O(b^{-3/2}),
 \qquad
 \int a_b^4g_b^2\,dy=O(b^{-2}).
\end{equation}
Because $g_b'=a_bg_b/2$ almost everywhere, the first estimate proves
the first part of \eqref{eq:CL-derivative-bounds}.
For $y>0$,
\[
 |a_b'(y)|=\frac1{2b}\left|\frac{B'(z)}{R(z)}\right|=O(b^{-1}),
 \qquad
 g_b''=\left(\frac{a_b^2}{4}+\frac{a_b'}2\right)g_b.
\]
For $y<0$, $g_b''=g_b/4$ and its squared norm on that half-line is
$e^{-b}/16$.  These facts and \eqref{eq:CL-score-estimates} prove the
second part of \eqref{eq:CL-derivative-bounds}.
The function $g_b$ is continuous at zero.  Its derivative jumps by
$J_b=-e^{-b/2}/(2b)$, giving the asserted singular part of its second
distributional derivative.  In particular $g_b\in H^1(\mathbb R)$;
we do not claim $g_b\in H^2(\mathbb R)$.

Let $v_b=(g_b'')_{\rm ac}$.  Integrating the distributional derivative
of $g_b'$ over an interval gives, for each real $s$,
\[
 \norm{T_sg_b'-g_b'}_2
 \leq |s|\norm{v_b}_2+|J_b|\sqrt{|s|}.
\]
For fixed $h=\log d$ and the usual orientation when $h<0$,
\[
 g_b-T_hg_b=\int_0^h T_sg_b'\,ds,
 \quad
 \norm{g_b-T_hg_b-hg_b'}_2
 \leq\frac{h^2}{2}\norm{v_b}_2
      +\frac23|h|^{3/2}|J_b|=O_h(b^{-1}).
\]
Since $\norm{g_b'}_2=O(b^{-1/2})$, squaring the last approximation
and applying \eqref{eq:exact-rank-translation} proves
\eqref{eq:CL-law-rigorous}.
\end{proof}

\subsection{Counting the exact discrete levels}

Let
\[
 F_r(q)=\prod_{\ell\geq0}(1+q^{2\ell+1})^r
       =\sum_{n\geq0}a_r(n)q^n,
 \qquad A_r(n)=\sum_{k=0}^n a_r(k),\qquad c=c_r.
\]
The coefficient $a_r(n)$ is the exact degeneracy of the level
$\varepsilon n$.  We take $A_r(-1)=0$.

\begin{lemma}[Counting with the required remainder]
\label{lem:tower-counting}
For each fixed positive integer $r$, as $n\to\infty$ through integers,
\begin{align}
 a_r(n)&=\frac{c^{1/4}}{2\sqrt\pi}\,n^{-3/4}
          e^{2\sqrt{cn}}\bigl(1+O_r(n^{-1/2})\bigr),
 \label{eq:tower-degeneracies}\\
 A_r(n)&=\frac{c^{-1/4}}{2\sqrt\pi}\,n^{-1/4}
          e^{2\sqrt{cn}}\bigl(1+O_r(n^{-1/2})\bigr)
        =I_0(2\sqrt{cn})\bigl(1+O_r(n^{-1/2})\bigr).
 \label{eq:tower-counting}
\end{align}
Moreover, if $L_r(t)=\log F_r(e^{-t})$, then
\begin{equation}\label{eq:tower-partition-estimates}
 L_r(t)=c/t+O_r(t),\qquad
 L_r'(t)=-c/t^2+O_r(1)
 \quad(t\downarrow0).
\end{equation}
\end{lemma}

\begin{proof}
Here is a direct saddle-point proof, including the minor-arc estimate
needed to exclude an unaccounted lattice periodicity.
Set $f(x)=\log(1+e^{-x})$.  Midpoint Euler--Maclaurin, applied to
$f$ and to $x^k f^{(k)}(x)$, gives for each fixed $k\geq0$
\begin{equation}\label{eq:partition-derivatives}
 L_r^{(k)}(t)=(-1)^k k!c\,t^{-k-1}+O_{r,k}(t^{1-k}).
\end{equation}
Indeed all the functions in question are smooth at zero, have
integrable second derivatives, and decay exponentially with their
derivatives at infinity; the midpoint-sum error is $O(t)$ before
multiplication by $t^{-k}$.  Integration by parts gives
$\int_0^\infty x^k f^{(k)}(x)\,dx=(-1)^k k!\int_0^\infty f(x)\,dx$,
and the last integral is $\pi^2/12$.
In a complex disk $|w-t|\leq t/2$ we also have
$|L_r^{(k)}(w)|\leq C_{r,k}t^{-k-1}$.
This follows either from the same sum or from Cauchy's estimate after
noting $|L_r(w)|\leq C_r/t$ for $\operatorname{Re}w\geq t/2$ and
$|\operatorname{Im}w|\leq t/2$; the factors are nonzero there.

Under the Boltzmann distribution with parameter $t$, the integer
energy $N_t$ is the sum of independent Bernoulli occupations times
their odd level numbers.  Its characteristic function is
\[
 \phi_t(\theta)=F_r(e^{-t+i\theta})/F_r(e^{-t}) .
\]
Retaining only odd levels $2j+1\leq1/t$ in the modulus of this product
gives
\[
 |\phi_t(\theta)|\leq
 \exp\!\left[-c_0\sum_{j=0}^{M-1}
              \sin^2((2j+1)\theta/2)\right],\qquad M\asymp t^{-1},
\]
with a positive $c_0$ independent of small $t$.  For $|\theta|\leq\pi$,
\begin{equation}\label{eq:odd-lattice-bound}
 \sum_{j=0}^{M-1}\sin^2((2j+1)\theta/2)
 \geq c_1 M\min\{1,M^2\theta^2\}.
\end{equation}
To check this, use $\sin u\geq2u/\pi$ when
$|\theta|\leq\pi/(2M)$.  In the middle region use the identity
$M/2-\sin(2M\theta)/(4\sin\theta)$ and
$|\sin\theta|\geq1/M$.  In the remaining region near $\pm\pi$,
write $\delta=\pi-|\theta|\leq\pi/(2M)$; the same expression is
$M/2+\sin(2M\delta)/(4\sin\delta)\geq M/2$.
Consequently
\begin{equation}\label{eq:characteristic-tail}
 |\phi_t(\theta)|\leq
       \exp[-c_2\min\{t^{-1},\theta^2t^{-3}\}].
\end{equation}

Take $t=\sqrt{c/n}$.  Fourier inversion gives
\[
 a_r(n)=\frac{e^{tn+L_r(t)}}{2\pi}
    \int_{-\pi}^{\pi}\phi_t(\theta)e^{-in\theta}\,d\theta.
\]
On the central arc set $\theta=t^{3/2}u$ with
$|u|\leq t^{-1/12}$.  Taylor's theorem and
\eqref{eq:partition-derivatives} give
\[
 L_r(t-i\theta)-L_r(t)-in\theta
 =-cu^2-ic\sqrt t\,u^3
   +O_r\!\left(tu^4+t^2u^2+t^{3/2}|u|
                    +t^{3/2}|u|^5\right).
\]
The error estimate is uniform on that arc.  The term of order
$\sqrt t$ is purely imaginary and odd.  Expanding the exponential,
integrating on the symmetric interval and using Gaussian domination
therefore gives
\[
 \int_{|\theta|\leq t^{3/2-1/12}}
          \phi_t(\theta)e^{-in\theta}\,d\theta
 =t^{3/2}\sqrt{\pi/c}\,(1+O_r(t)).
\]
For clarity, the integrated error is bounded by a constant times
$t\int_{\mathbb R}e^{-c u^2/2}(1+|u|^{10})\,du$;
odd imaginary terms integrate to zero because the values at $u$ and
$-u$ are complex conjugates.  The complement of the central arc is
$O(\exp[-c_3t^{-1/6}])$ by \eqref{eq:characteristic-tail}.
Finally $tn+L_r(t)=2\sqrt{cn}+O_r(t)$, which proves
\eqref{eq:tower-degeneracies} with the asserted relative $O(t)$
remainder, rather than merely an ordinary central limit estimate.

To sum the coefficients, put $H(x)=I_0(2\sqrt{cx})$ for $x\geq0$.
The fixed-order Bessel expansion implies
\[
 H'(x)=\sqrt{c/x}\,I_1(2\sqrt{cx})
 =\frac{c^{1/4}}{2\sqrt\pi}\,x^{-3/4}e^{2\sqrt{cx}}
      (1+O_r(x^{-1/2})).
\]
Thus $a_r(k)=H'(k)(1+O_r(k^{-1/2}))$ for all sufficiently large $k$.
The function $H'$ is eventually increasing.  Integral comparison
then gives $\sum_{k\leq n}H'(k)=H(n)+O_r(H'(n)+1)$.
The total of the relative-error terms is $O_r(n^{-1/2}H(n))$:
split at $n/2$, use the exponential separation below $n/2$, and
bound $k^{-1/2}\leq\sqrt{2/n}$ above it.  Since
$H'(n)/H(n)=O_r(n^{-1/2})$, this proves \eqref{eq:tower-counting}.
\end{proof}

\begin{remark}[Independent number-theoretic check]
For $r=1,2$ the coefficient estimate also follows from Chern's eta-quotient theorem~\cite[Theorem~1.1]{Chern19}.  Indeed
\[
 F_r(q)=\frac{(q^2;q^2)_\infty^{2r}}
                 {(q;q)_\infty^r(q^4;q^4)_\infty^r}.
\]
In his notation $m=(1,2,4)$, $\delta=(-r,2r,-r)$,
$\Delta_1=0$, $\Delta_2=-r$, and
$\Delta_3(k)$ equals $r/4,-2r,r$ when $k$ is odd,
$k\equiv2\pmod4$, and $k\equiv0\pmod4$, respectively.
His hypothesis (1.8) holds for $1\leq r\leq24$.
The $k=1$ term is
$\sqrt{c/(n-r/24)}I_1(2\sqrt{c(n-r/24)})$; every other exponentially
growing term has at most half its exponential rate.  This checks the
normalization of \eqref{eq:tower-degeneracies}.  The direct proof above
does not need the restriction $r\leq24$.
\end{remark}

\subsection{Uniform inversion and completion of the proof}

\begin{lemma}[Uniform microscopic comparison]
\label{lem:tower-uniform-comparison}
Let $p_{r,\varepsilon}(x)=p_j$ on $(j-1,j]$ for the ordered spectrum
\eqref{eq:tower-spectrum-rigorous}, and put $b=c_r/\varepsilon$.
There is a constant $C_r<\infty$ such that, for all sufficiently
small $\varepsilon>0$,
\begin{equation}\label{eq:tower-uniform-logratio}
 \underset{x>0}{\operatorname{ess\,sup}}
 \left|\log\frac{p_{r,\varepsilon}(x)}{q_b(x)}\right|
 \leq C_r\varepsilon,
 \qquad
 \norm{\sqrt{p_{r,\varepsilon}}-\sqrt{q_b}}_2=O_r(\varepsilon).
\end{equation}
\end{lemma}

\begin{proof}
Fix $r$ and write $c=c_r$ and $H(u)=I_0(2\sqrt{cu})$.
For $x>1$ define $u(x)=H^{-1}(x)$, and put $u(x)=0$ for
$0<x\leq1$.  If $A_r(n-1)<x\leq A_r(n)$, the exact level at rank
$x$ is $n$.  We claim
\begin{equation}\label{eq:bounded-inverse-energy}
 |u(x)-n|\leq K_r
\end{equation}
with a constant independent of both $n$ and $x$.
By \eqref{eq:tower-counting},
$\log A_r(n)-\log H(n)=O_r(n^{-1/2})$, whereas
\[
 \frac{d}{du}\log H(u)=\sqrt{c/u}\,
               \frac{I_1(2\sqrt{cu})}{I_0(2\sqrt{cu})}
             =\sqrt{c/u}+O_r(u^{-1}).
\]
Choose a fixed, sufficiently large $K_r$.  For all sufficiently large
$n$, these two estimates imply
$H(n-K_r)\leq A_r(n-1)$ and $A_r(n)\leq H(n+K_r)$.
Monotonicity of $H$ proves \eqref{eq:bounded-inverse-energy} there.
There are only finitely many smaller levels, and their rank intervals
are bounded, so increasing $K_r$ proves it everywhere.
Empty intervals caused by zero degeneracies require no treatment;
for example $a_1(2)=0$, while $a_1(n)>0$ for every $n\geq3$.

At the rank under consideration,
\[
 p_{r,\varepsilon}(x)=e^{-L_r(\varepsilon)-\varepsilon n},
 \qquad q_b(x)=e^{-b-\varepsilon u(x)}.
\]
Since $L_r(\varepsilon)=b+O_r(\varepsilon)$, the logarithm of their
ratio is uniformly $O_r(\varepsilon)$, proving the first claim.
Both functions are normalized, and pointwise
$|\sqrt p-\sqrt q|\leq(e^{C_r\varepsilon/2}-1)\sqrt q$ after
enlarging the constant, so integration proves the second claim.
\end{proof}

\begin{proof}[Proof of Theorem~\ref{thm:exact-tower-law}]
The entropy of the odd tower is
$S=L_r(\varepsilon)-\varepsilon L_r'(\varepsilon)$. Equation
\eqref{eq:partition-derivatives} gives
\[
 S=2c_r/\varepsilon+O_r(\varepsilon),\qquad
 V=\varepsilon^2 L_r''(\varepsilon)
   =2c_r/\varepsilon+O_r(\varepsilon).
\]
Apply Proposition~\ref{prop:xx-count-to-error} with $W=\varepsilon^{-1}$
and $L(z)=L_r(\varepsilon z)$. The real and complex derivative bounds
in Lemma~\ref{lem:tower-counting} give $v_\beta\asymp_r W$ and the
required analytic bounds. The characteristic function of the physical
energy under tilt $\beta$ is $\phi_{\varepsilon\beta}(\varepsilon t)$
in the notation of that lemma. For $|t|\le\pi W$, the odd-lattice
estimate~\eqref{eq:characteristic-tail} gives
\[
 |\phi_{\varepsilon\beta}(\varepsilon t)|
 \le \exp\bigl(-c W\min\{1,t^2\}\bigr)
\]
uniformly in $\beta$ near one. This proves the hypotheses of the
counting criterion. Its conclusion is
\[
 \mathcal E_d=\frac{(\log d)^2}{4V}+O_{r,d}(\varepsilon^2)
      =\frac{(\log d)^2\varepsilon}{8c_r}+O_{r,d}(\varepsilon^2),
 \qquad
 S\mathcal E_d=\frac{(\log d)^2}{4}+O_{r,d}(\varepsilon).
\]
\end{proof}


\subsection{A transfer condition for other spectra}

The same leading law transfers to other spectra under a quantitative Hellinger estimate.

\begin{corollary}[A sufficient transfer condition]
\label{cor:CL-physical-transfer}
Let $(p_{N,j})$ be any sequence of decreasing probability spectra, let
$p_N(x)=p_{N,j}$ on $(j-1,j]$, and let $b_N\to\infty$.  Suppose
\begin{equation}\label{eq:physical-transfer-hypotheses}
 S(p_N)/(2b_N)\longrightarrow1,
 \qquad
 \norm{\sqrt{p_N}-\sqrt{q_{b_N}}}_2=o(b_N^{-1/2}).
\end{equation}
Then, for every fixed integer $d\geq2$,
$\mathcal E_d(p_N)S(p_N)\to(\log d)^2/4$.
In particular the second hypothesis holds if
$\operatorname{ess\,sup}_{x>0}|\log(p_N(x)/q_{b_N}(x))|
=o(b_N^{-1/2})$.
\end{corollary}

\begin{proof}
The first assertion follows directly from
Lemma~\ref{lem:hellinger-transfer} and
Proposition~\ref{prop:CL-sharp-law}.  The final sufficient criterion
follows from the same pointwise square-root estimate used in
Lemma~\ref{lem:tower-uniform-comparison}.
\end{proof}

For finite-dimensional physical spectra the final, global log-ratio
criterion cannot hold because the spectrum has a finite-rank tail;
the Hellinger condition in \eqref{eq:physical-transfer-hypotheses} is
the usable criterion and includes the tail mass.  It may also be
replaced by the strictly weaker, direct condition
$\norm{(I-T_{\log d})(g_{p_N}-g_{q_{b_N}})}_2=o(b_N^{-1/2})$.
Neither leading R\'enyi entropy asymptotics nor weak Gaussian
convergence of logarithmic rank proves these microscopic conditions.
Such a transfer therefore requires more than the leading entropy asymptotic.
